\section{Evaluation Protocol}
\label{sec:study-design}

The protocol recombines independent submissions while fixing the image, task condition and target reference. It separates the effect of adding people from changes in profile composition, individual membership and retained background (Table~\ref{tab:principal-analysis-sets}). Cross-view and doorway comparisons then examine changes in visual evidence and enclosure choice.

\subsection{Task, participants and reference targets}

Participants marked the walls enclosing the camera, including architectural recesses, while ignoring movable furniture and decoration. The main task assumes vertical walls along two orthogonal plan-view axes, a closed single-level footprint, and single floor and ceiling planes. Openings are closed at identifiable structural divisions such as doorframes or paired wall ends. Doorway-boundary viewpoints place the camera beneath a doorframe and expose adjoining enclosures. These cases require an explicit target and are analyzed separately from ordinary rooms; oblique walls, sloped ceilings and split levels form an out-of-scope (OOS) cohort.

Participants were recruited from university cohorts and Upwork, received an average of USD~0.50 per panorama, and consented to duration recording and research use. After instruction and practice, \textit{Manual} annotators placed and ordered upper--lower pairs from scratch using an interactive 3D preview. \textit{Semi} annotators edited HoHoNet predictions \cite{hohonet2021} or six controlled-defect initializations on separate panoramas. The main collection contains 6,512 Manual submissions from 60 people on 206 panoramas in 22 buildings and 960 Semi submissions from 59 people on 37 panoramas in 12 buildings; 57 of the 62 total participants contributed to both modes (Supplements B1, C12, C20).

Before the study, three experts with experience annotating 100 panoramas and detailed knowledge of Manhattan layouts reviewed MatterportLayout references \cite{chang2017matterport3d,zou2021manhattan}, resolved target rules, corrected geometric errors, and catalogued natural Semi initialization defects. Reference corner counts, self-occlusion and expert inspection define easy, moderate and difficult tiers. Experts screened doorway responses for insufficient points or uninterpretable enclosures while blinded to $Q$. The fixed 17-image screening comparison retains the same whole-space reference before and after screening; separate target-assigned pools address enclosure choice (Supplements C15, C17, C20).

\subsection{Controlled team comparisons}

For RQ2, Monte Carlo sampling with 256 uniformly random entry orders per eligible pool form nested teams of 3--30 people on a fixed panel of 171 panoramas. For RQ3, profiles are derived from each person's history outside the target building. Structural classes C1--C3 represent simplifying, typical and turn-retaining tendencies. Their features cover corner counts, short edges, near-collinear vertices and reference matches. Within each training fold, robust scaling, domain-wise PCA and domain balancing precede Ward clustering; classes are ordered by difficult-image turns and short edges, then overall turns. Effective turns exceed $5^\circ$; short edges occupy less than 4\% of the perimeter with endpoint turns of at least $15^\circ$. Consensus-response (R1--R2) profiles summarize historical effects on group geometry; joint profiles (J1--J4) combine structural, geometric, historical-quality and duration features (Supplements C6--C9).

\begin{table*}[t]
\centering\small
\caption{Principal controlled comparisons. The image and target are fixed except where viewpoint or target is the studied factor. Supports overlap; Supplement C20 reports full coverage.}
\label{tab:principal-analysis-sets}\label{tab:analysis-support}
\setlength{\tabcolsep}{5pt}\renewcommand{\arraystretch}{1.12}
\begin{tabular}{@{}>{\raggedright\arraybackslash}p{.19\textwidth}>{\raggedright\arraybackslash}p{.31\textwidth}>{\raggedright\arraybackslash}p{.44\textwidth}@{}}
\toprule
Comparison & Support & Controlled change and outcome\\
\midrule
Group size & 171 images; 22 buildings; 256 orders/image & Add people; quality distribution and completion\\
Shared structure & 70 Manual teams; 24 images; 14 buildings & Same submissions; floor coverage, turns and upper boundaries\\
Matched replacement & 69 images; 20 buildings; 146/147 scorable events & Replace 7/14 members; excess magnitude and direction\\
Retained background & 276 blocks; 7 images & Same exchange, different retained members; direction\\
Cross-view structure & 128 view pairs; 65 images; 13 buildings & Same and different people across room views; repetition and transfer\\
Doorway screening & 17 images; 9 buildings; 584/439 responses before/after & Screen responses against the same whole-space reference\\
\bottomrule
\end{tabular}
\end{table*}

A direct replacement compares $A=D\cup X$ with $C=D\cup Z$, retaining members $D$ while exchanging equal-sized blocks $X$ and $Z$. The matched A--B--C design adds $B=D\cup Y$: $X$ and $Y$ share a profile, $Z$ has a contrasting profile, and $D,X,Y,Z$ are mutually disjoint. It separates excess change magnitude from direction:
\begin{align}
G_S&=\frac{|Q_C-Q_A|+|Q_C-Q_B|}{2}-|Q_B-Q_A|,\label{eq:symmetric-excess}\\
\Delta_{CB}&=Q_C-Q_B.\label{eq:directional-change}
\end{align}
Effects are averaged over valid exchanges and backgrounds within each image, then equally over images. Fixed-size profile ratios and whole-profile mixtures test composition; holding $(X,Z)$ fixed while changing $D$ tests background dependence (Supplement C17).

\subsection{Method assessment and outcome measures}

The $Q$ form and scales were fixed using controlled geometric perturbations before human validation and team analysis. Measurements use 2,048 boundary and 4,096 height samples per direction. Consensus uses $g=7.5^\circ$ per skipped observation to accommodate within-corner variation while distinguishing nearby turns. Supplement C16 records these settings and the existing score-scale sensitivity checks.

For human validation, two Manual annotations were randomly sampled from each of 30 panoramas in 17 buildings: nine easy, nine moderate, nine difficult and three doorway images. Eight randomly invited annotators graded them as good, medium or poor and marked visible defects using randomized displays of panorama overlays, shared-scale footprints and 3D views. Shared-team assessment compares the same teams using paired-corner consensus and BEV region voting, which retains planar cells covered by at least half the team \cite{lee2018}.

Team outcomes are mean $Q$, within-image P10, between-team standard deviation, and complete-and-evaluable rate. Quality comparisons use scorable outputs; completion counts all attempts. For group-size curves, 2,000 bootstrap resamples of entry trajectories estimate Monte Carlo precision, while 2,000 building resamples describe variation across the 22 buildings with the observed annotators retained.

Offline selection and stopping use recorded annotation person-time. Quality-and-time selection exchanges members of a historical-$Q$ team for faster people within a two-point historical-quality tolerance. Source-history stopping uses other views to select the earliest team size within two points of their expected peak, continuing to obtain a complete polygon. It is compared with fixed seven-person teams and a three-step corner-count stability rule; source-history acquisition is included in project labor (Section~\ref{sec:labor-history}, Supplements C8--C9, C20).
